\section{Kokeet}
Ensimmäisessä kokeessa tarkastellaan Cortex Analystin konsistenssia, eli sitä, tuottaako järjestelmä saman SQL-kyselyn, kun sama luonnollisen kielen kysymys esitetään useita kertoja.

Toisessa kokeessa tarkastellaan järjestelmän herkkyyttä sellaisille syötteen muutoksille, joiden ei pitäisi muuttaa kysymyksen merkitystä. Näissä kokeissa muodostetaan useita eri sanamuotoja samasta tiedontarpeesta ja arvioidaan, tuottaako järjestelmä niistä olennaisesti saman SQL-kyselyn.

\paragraph{Konsistenssikoe}
Koska kielimallit ovat tilastollisia järjestelmiä, sama syöte ei välttämättä aina johda täsmälleen samaan tulosteeseen. Tietokantakyselyiden tapauksessa tämä on olennainen kysymys, koska käyttäjät odottavat järjestelmältä ennustettavaa ja toistettavaa toimintaa.

Konsistenssikokeessa sama luonnollisen kielen kysymys esitetään järjestelmälle 10 kertaa. Jokaisesta suorituksesta tallennetaan järjestelmän tuottama SQL-kysely, joita verrataan keskenään. Käytetyt kysymykset ovat seuraavat:
\begin{enumerate}[label=\Alph*.]
    \label{it:q_consistency}
  \item \textit{How many total contacts did we handle in January 2026?}
  \item \textit{What's our answer rate for calls in February 2026?}
  \item \textit{What percentage of contacts that were answered within our service level target were closed by the customer versus closed by the agent?}
  \item \textit{What was our peak contact hours by day of week? Show me the distribution of when contacts arrive throughout the day, broken down by hour and weekday.}
  \item \textit{What's the success rate of our callbacks? Show me how many callbacks resulted in a conversation versus no answer, and what the average time between the original inbound contact and the callback was.}
\end{enumerate}

Määritellyssä semanttisessa näkymässä on kuvattu metriikoihin SQL-kyselyt, jotka vastaavat suoraan kysymyksiin A ja B. Kysymykset C-E ovat huomattavasti monimutkaisempia, ja niihin ei ole ennalta validoituja SQL-kyselyitä semanttisessa näkymässä. Kysymys D vaatii, että kentästä \texttt{CONTACT\_CONNECTED\_DATETIME} erotetaan kellonaika tunnin tarkkuudella, ja viikonpäivä pitää hakea kalenteritaulusta. Kysymys E on rakenteeltaan monitulkintainen. Erityisesti takaisinsoiton ja alkuperäisen sisääntulevan kontaktin välisen ajan keskiarvoa koskeva osa voidaan tulkita vaativan kahden eri rivin välistä vertailua, vaikka käytetyssä aineistossa vastaus voidaan muodostaa yhdeltä takaisinsoittoa kuvaavalta riviltä.

\paragraph{Sensitiivisyyskoe}
Kokeen tavoitteena on arvioida, tulkitseeko Cortex Analyst semanttisesti saman sisältöiset kysymykset johdonmukaisesti vai johtavatko kielelliset erot erilaisiin SQL-kyselyihin. Jos järjestelmä on hyvin herkkä muotoilun vaihtelulle, luonnollisen kielen käyttö rajapintana muuttuu käyttäjän kannalta epävarmaksi. Jos taas järjestelmä tuottaa olennaisesti saman kyselyn eri sanamuodoista huolimatta, tämä tukee sen käytettävyyttä käytännön työssä.

Koe toteutetaan muodostamalla kustakin konsistenssikokeessa käytetystä kysymyksestä 4 vaihtoehtoista versiota, joissa pyritään säilyttämään kysymyksen merkitys, mutta muokkaamaan sanamuotoa. Järjestelmän tuottamat SQL-kyselyt tallennetaan ja vertaillaan keskenään. Käytetyt kysymykset ovat seuraavat:
\begin{enumerate}[label=\Alph*.]
    \label{it:q_sensitivity}
  \item Variaatioita kysymyksestä A:
    \begin{enumerate}[label=\arabic*.]
      \item \textit{What was the total number of contacts we handled in January 2026?}
      \item \textit{How many contacts did we handle overall in January 2026?}
      \item \textit{What is the total contact volume for January 2026?}
      \item \textit{How many total customer contacts were handled during January 2026?}
    \end{enumerate}
  \item Variaatioita kysymyksestä B:
    \begin{enumerate}[label=\arabic*.]
      \item \textit{What was our call answer rate in February 2026?}
      \item \textit{What percentage of incoming calls were answered in February 2026?}
      \item \textit{What percentage of calls were answered in February 2026?}
      \item \textit{How did our call answer rate perform during February 2026?}
    \end{enumerate}
  \item Variaatioita kysymyksestä C:
    \begin{enumerate}[label=\arabic*.]
      \item \textit{Of the contacts answered within our service-level target, what share were closed by the customer compared to those closed by the agent?}
      \item \textit{For interactions resolved within the service level target, what percentage were customer-closed versus agent-closed?}
      \item \textit{Within our service level target, how is the closure split: what percent of contacts were closed by customers versus agents?}
      \item \textit{Among contacts answered within SLA, what proportion ended with customer closure versus agent closure?}
    \end{enumerate}
  \item Variaatioita kysymyksestä D:
    \begin{enumerate}[label=\arabic*.]
      \item \textit{When do contacts most frequently arrive during the week? Break this down by weekday and hour.}
      \item \textit{Show the hourly contact-arrival distribution for every day of the week, highlighting which hours were highest on each weekday.}
      \item \textit{What times during the day (hour-by-hour) drive the most contacts, split by weekday. Please summarize peak hours and the overall distribution.}
      \item \textit{Across the week, when do contacts arrive most often? Provide an hour-by-hour breakdown for each weekday and identify the peak contact hours.}
    \end{enumerate}
  \item Variaatioita kysymyksestä E:
    \begin{enumerate}[label=\arabic*.]
      \item \textit{What percentage of our callback attempts lead to an actual conversation? Please break down callbacks that reached a conversation vs. those with no answer, and calculate the average delay from the original inbound to the callback.}
      \item \textit{How effective are our callbacks? Show the success rate by counting callbacks that resulted in a conversation versus no answer, and report the average time-to-callback from the initial inbound contact.}
      \item \textit{For our callback campaign, what's the outcomes split and average response speed? Provide counts of conversations vs. no-answer callbacks, the resulting success rate, and the average elapsed time between inbound contact and callback.}
      \item \textit{Can you analyze callback performance for me: success rate (conversation vs no answer) with the number of callbacks in each outcome category, plus the average time lag from the inbound contact to when the callback occurred?}
    \end{enumerate}
\end{enumerate}

Jokainen kysymys molemmissa kokeissa esitetään järjestelmälle 10 kertaa ja jokainen yksittäinen suoritus tehdään tyhjällä sessiolla. Testit suoritettiin 23.03.2026-28.03.2026 välisenä aikana käyttäen samaa semanttista mallia.

Kymmenen toistoa valittiin käytännöllisenä kompromissina, joka mahdollistaa vaihtelun havaitsemisen ilman kohtuuttoman suurta manuaalisen arvioinnin työmäärää.
